% =============================================================================
\section{Patient Input Pipeline}
% =============================================================================

\begin{frame}{Step 1: How data enters the system}
  \textbf{Session start:} patient describes symptoms via the Curatio chatbot.

  \vspace{0.3em}
  \begin{columns}[T]
    \column{0.48\textwidth}
    \begin{block}{Two input modes}
      \begin{itemize}
        \item \textbf{Typed text} --- stored directly as raw complaint
        \item \textbf{Voice note} --- Twi / local Ghanaian language
      \end{itemize}
    \end{block}

    \column{0.48\textwidth}
    \begin{block}{5-step pipeline}
      \footnotesize
      \begin{enumerate}
        \item Patient input (this slide)
        \item Text example
        \item Speech \& translation
        \item Fine-tuned BioBERT
        \item Fusion $\to$ colour $C$
      \end{enumerate}
      \normalsize
      TEWS and Bayesian join at Step 5.
    \end{block}
  \end{columns}

  \plain{Whether the patient types or speaks, the goal is the same: produce English text $X$ for the NLP layer.}
\end{frame}

\begin{frame}{Step 2: Example text input}
  \textbf{Typed complaint} is stored as \texttt{chief\_complaint\_raw} and becomes text $X$ (max 128 tokens).

  \examplebox{
    \textbf{Patient ID:} TG-UXRGA9UCO\\
    \textbf{Complaint:} \textit{``thunderclap headache, worsening with movement''}\\
    \textbf{Nurse ground truth:} Acuity Level 2 --- \sats{triageOrange}{Orange} (highly urgent)
  }

  \vspace{0.3em}
  Voice input skips this direct path and goes to Step 3 (ASR + translation) before BioBERT.

  \plain{Real hospital phrasing --- not textbook English --- is exactly what fine-tuning teaches BioBERT to handle.}
\end{frame}

\begin{frame}{Step 3: Speech and translation}
  \begin{equation*}
    X = g_{\mathrm{trans}}\!\bigl(g_{\mathrm{ASR}}(\mathrm{audio})\bigr)
  \end{equation*}

  \begin{itemize}
    \item $g_{\mathrm{ASR}}$ --- automatic speech recognition (Twi $\to$ transcript)
    \item $g_{\mathrm{trans}}$ --- medical translation (Twi $\to$ English)
    \item Typed English from Step 2 bypasses both functions
  \end{itemize}

  \examplebox{
    \textbf{Patient TG-001} (Twi voice):
    \begin{enumerate}
      \item Audio: \textit{``Me koma mu yɛ me ya, me home yɛ me''}
      \item ASR: \textit{Me koma mu yɛ me ya}
      \item Translation: \textit{``I have chest pain, I feel unwell''}
      \item $X \to$ BioBERT in Step 4
    \end{enumerate}
  }

  \plain{$g_{\mathrm{ASR}}$ = ears writing words; $g_{\mathrm{trans}}$ = medical interpreter. BioBERT reads \textbf{English only}.}
\end{frame}

\begin{frame}{Step 3: Pipeline diagram}
  \begin{center}
    \resizebox{0.85\textwidth}{!}{%
    \begin{tikzpicture}[node distance=0.8cm]
      \node[flowbox, minimum width=2cm] (audio) {Voice\\audio};
      \node[flowbox, right=of audio] (asr) {$g_{\mathrm{ASR}}$\\Transcribe};
      \node[flowbox, right=of asr] (trans) {$g_{\mathrm{trans}}$\\Translate};
      \node[flowbox, right=of trans, fill=triageGreen!10] (x) {English $X$};
      \node[flowbox, below=1cm of audio, minimum width=2cm] (text) {Typed\\text};
      \draw[flowarr] (audio) -- (asr) -- (trans) -- (x);
      \draw[flowarr] (text) -| (x);
    \end{tikzpicture}}
  \end{center}

  \vspace{0.5em}
  Both paths converge on the same $X$ before the BioBERT layer.
\end{frame}
